\documentclass[11pt,a4paper]{article}
\usepackage{brainteaser}

\begin{document}
\brainteasertitle{10}{Deux six consécutifs}

\begin{statementbox}
On lance un dé équilibré à six faces jusqu'à obtenir deux $6$ consécutifs.
Quel est le nombre moyen de lancers nécessaires?
\end{statementbox}

\begin{hintsbox}
\begin{enumerate}
  \item L'historique complet est-il utile, ou suffit-il de retenir une
  information sur le dernier lancer?
  \item Définir $E_0$, l'espérance restante lorsque le dernier lancer n'est
  pas un $6$, et $E_1$, celle obtenue lorsque le dernier lancer est un $6$.
  \item Conditionner chaque espérance sur le prochain lancer. Ne pas oublier
  que ce lancer compte toujours pour $1$.
  \item Après un échec depuis l'état $E_1$, dans quel état revient-on?
\end{enumerate}
\end{hintsbox}

\correction

Depuis l'état $E_0$, un $6$ fait passer dans l'état $E_1$, tandis qu'un autre
résultat laisse dans $E_0$:
\[
E_0=1+\frac16E_1+\frac56E_0.
\]
Depuis l'état $E_1$, un nouveau $6$ termine le processus; un autre résultat
ramène à $E_0$:
\[
E_1=1+\frac56E_0.
\]

La première équation donne $E_0=6+E_1$. En remplaçant $E_1$,
\[
E_0=7+\frac56E_0,
\]
d'où
\[
E_0=42,
\qquad
E_1=36.
\]

Au début, aucun $6$ précédent n'a été obtenu: on part de $E_0$.

\begin{resultbox}
\[
\boxed{42\text{ lancers en moyenne}.}
\]
\end{resultbox}

\end{document}
